% Generated draft: requires mathematical and source review.
\documentclass[11pt]{article}
\usepackage[margin=25mm]{geometry}
\usepackage{fontspec}
\setmainfont{Latin Modern Roman}
\usepackage{amsmath,amssymb,mathtools,unicode-math}
\setmathfont{Latin Modern Math}
\usepackage{enumitem,needspace,xurl,hyperref}
\setlength{\parindent}{0pt}
\setlength{\parskip}{5pt}
\setlength{\emergencystretch}{3em}
\newcommand{\N}{\mathbb{N}}
\newcommand{\Z}{\mathbb{Z}}
\newcommand{\Q}{\mathbb{Q}}
\newcommand{\R}{\mathbb{R}}
\newcommand{\C}{\mathbb{C}}
\newcommand{\F}{\mathbb{F}}
\newcommand{\A}{\mathbb{A}}
\newcommand{\PP}{\mathbb{P}}
\newcommand{\E}{\mathbb{E}}
\newcommand{\eps}{\varepsilon}
\newcommand{\dd}{\,\mathrm{d}}
\begin{document}
\section*{UnsolvedMath ID: 801196}
\textbf{Source problem code:} catalogue entry
\subsection{Definitions and mathematical statement}
Let \(P=(V,\leq_P)\) be a finite poset and let \(n=|V|\). Say that \(P\) is \((3+1)\)-free if there are no four distinct elements \(a,b,c,d\in V\) such that \(a<_P b<_P c\) and \(d\) is incomparable with each of \(a,b,c\). Thus the forbidden configuration is an induced subposet consisting of a three-element chain and an isolated element. The incomparability graph \(G=\operatorname{inc}(P)\) has vertex set \(V\), with an edge between distinct vertices exactly when they are incomparable in \(P\).

A proper coloring of \(G\) is a map \(\kappa:V\to\{1,2,3,\ldots\}\) satisfying \(\kappa(u)\ne\kappa(v)\) for every edge \(\{u,v\}\). For commuting indeterminates \(x=(x_1,x_2,\ldots)\), define the chromatic symmetric function by
\[
X_G(x)=\sum_{\kappa\text{ proper}}\prod_{v\in V}x_{\kappa(v)}
=\sum_{\kappa\text{ proper}}\prod_{i\geq1}x_i^{a_i(\kappa)},
\qquad a_i(\kappa)=|\kappa^{-1}(i)|.
\]
For \(r\geq1\), let
\[
e_r(x)=\sum_{1\leq i_1<\cdots<i_r}x_{i_1}\cdots x_{i_r},
\qquad e_0(x)=1.
\]
For a partition \(\lambda=(\lambda_1,\ldots,\lambda_\ell)\) of \(n\), set \(e_\lambda=\prod_{j=1}^{\ell}e_{\lambda_j}\). When \(n=0\), use the empty partition and \(e_{\varnothing}=1\).

\textbf{Conjecture (Stanley--Stembridge; Stanley's Problem 21).} For every finite \((3+1)\)-free poset \(P\), the unique elementary-basis expansion
\[
X_{\operatorname{inc}(P)}(x)=\sum_{\lambda\vdash n}c_\lambda(P)e_\lambda(x)
\]
satisfies \(c_\lambda(P)\in\Z_{\geq0}\) for every partition \(\lambda\vdash n\). This property is called \(e\)-positivity. The supplied registry labels the problem open; that label is not an independently verified assessment of its current mathematical status.
\subsection{Short English statement}
The Stanley--Stembridge conjecture asks whether the chromatic symmetric function of the incomparability graph of every finite \((3+1)\)-free poset has nonnegative coefficients when expanded in elementary symmetric functions.
\subsection{Sources}
\begin{itemize}
\item Local UnsolvedMath registry, permanent ID \(801196\), source code \texttt{catalogue entry}, in the AMR Open Problem Lists collection. The supplied record attributes the problem to Stanley, lists proposed year \(1999\), and labels its status \texttt{open}.
\item Richard P. Stanley, \emph{Positivity Problems and Conjectures in Algebraic Combinatorics} (1999), Problem 21. Supplied source URL: \url{https://math.mit.edu/~rstan/papers/problems.pdf}. The record reports access on \texttt{2026-10-16} and extraction marked \texttt{pdf+html\_verified}; these are supplied provenance claims, not independently verified here.
\item The supplied background attributes the conjecture to Richard P. Stanley and John R. Stembridge, \emph{J. Combin. Theory Ser. A} \textbf{62} (1993), 261--279, Conjecture 5.5, and to Stanley, \emph{Adv. Math.} \textbf{111} (1995), 166--194. Article titles and separate links for these references are not supplied.
\end{itemize}
\end{document}
